\documentclass[12pt]{article}
\usepackage[utf8]{inputenc}
\usepackage[T1]{fontenc}
\usepackage{hyperref}
\title{Taman Language Resources for K-12 Curriculum Development}
\author{Compiled by Claude Sonnet 5 (claude-sonnet-5) for Andr\'as Kornai}
\date{}
\begin{document}
\maketitle
\section{Introduction}
Taman (ISO 639-3: tcl; Glottocode tama1328) is a Sino-Tibetan language once spoken in Htamanthi village, Homalin Township, Sagaing Region, northern Myanmar. The sources I found disagree on when it fell out of use. Wikipedia gives 1931 and notes one person who remembered some of it in 2015. Joshua Project says it was functionally dormant by the 1930s and that the last known speaker died in the 1990s. Its place within Tibeto-Burman is uncertain. Benedict and Shafer put it in Luish, and Huziwara argues it cannot be securely assigned there. Do not confuse it with the ASJP entry ``Taman Myanmar'', which carries a different ISO code and Glottocode.
\section{Primary Educational Resources}
I found no primers, textbooks or school materials for Taman. It is not known to be taught in schools.
\section{Linguistic and Reference Materials}
\begin{itemize}
\item Wikipedia, Taman language (Myanmar): summary of the sources, including a 75-word list compiled by Brown in 1911 and a 2016 fieldwork note by Keisuke Huziwara. \url{https://en.wikipedia.org/wiki/Taman_language_(Myanmar)}
\item Glottolog languoid page for Taman (tama1328). \url{https://glottolog.org/resource/languoid/id/tama1328.bigmap.html}
\item Endangered Languages Project entry. \url{https://endangeredlanguages.com/lang/4274/guide?hl=de}
\end{itemize}
\section{Religious and Community Resources}
\begin{itemize}
\item Joshua Project, Taman in Myanmar (people-group profile). \url{https://joshuaproject.net/people_groups/15224/BM}
\item Joshua Project language page for tcl. \url{https://joshuaproject.net/languages/tcl}
\end{itemize}
I found no scripture translations or community organisations.
\section{Digital and Audio Resources}
I found no audio, video or corpus resources in the search results.
\section{Curriculum Development Context}
The only substantial record found is the 1911 word list of about 75 items, plus phrases and a short song recalled in 2016 by one woman who was not fluent. A curriculum that uses the language itself is not feasible from these sources. Any heritage or revitalization work would have to start from the documentation by Huziwara and the 1911 list, ideally with the descendant community and linguists. I did not locate the primary publications in search.
\end{document}